Fix an integer $k\ge 2$ and $n\ge 1$. Let $\mathsf{SAW}_{k}(n)$ be the set of all length-$n$ nearest-neighbour walks
\[
\omega=(\omega_0,\omega_1,\dots,\omega_n),\qquad \omega_0=0\in\mathbb Z^k,\qquad \|\omega_{t+1}-\omega_t\|_1=1,
\]
that are \emph{self-avoiding}: $\omega_i\neq \omega_j$ for all $0\le i<j\le n$.
Let $\mathbb P_{k,n}$ be the uniform probability measure on $\mathsf{SAW}_{k}(n)$, and let $X_n(\omega):=\omega_n$ be the endpoint.
Let $\|\cdot\|$ denote the Euclidean norm on $\mathbb R^k$, and define
\[
d_k(n):=\mathbb E_{k,n}\big[\|X_n\|\big].
\]
The questions are:

(1) Is it true that $\displaystyle \lim_{n\to\infty}\frac{d_2(n)}{\sqrt n}=+\infty$?

(2) Is it true that for each fixed $k\ge 3$ there is a constant $C_k$ such that $d_k(n)\le C_k\sqrt n$ for all $n\ge 1$?

\emph{Ambiguities in the literal statement.} The phrase ``expected distance'' is not explicit about the metric; I fix Euclidean norm because it is standard in the self-avoiding walk literature and matches the scaling statements in the problem text. Also ``taking $n$ random steps ... conditional on no self intersections'' is interpreted as conditioning simple random walk on self-avoidance; Lemma~\ref{lem:529-uniform} shows this is equivalent to the uniform measure on $\mathsf{SAW}_k(n)$.
